\documentclass[12pt]{amsart}
\usepackage{amsmath, amssymb, amsthm}
\usepackage[margin=1in]{geometry}
\usepackage{graphicx}
\usepackage{hyperref}
\usepackage{xcolor}
\usepackage{booktabs}
\usepackage{microtype}
\hypersetup{
  colorlinks=true,
  linkcolor=black,
  citecolor=black,
  urlcolor=black
}
\pagestyle{plain}

\newcommand{\derived}{\textsc{[derived]}}
\newcommand{\openq}{\textsc{[open]}}
\newcommand{\assumed}{\textsc{[assumed]}}
\newcommand{\conceptual}{\textsc{[conceptual]}}
\newcommand{\interpretive}{\textsc{[interpretive]}}

\newtheorem{theorem}{Theorem}[section]
\newtheorem{lemma}[theorem]{Lemma}
\newtheorem{proposition}[theorem]{Proposition}
\newtheorem{corollary}[theorem]{Corollary}
\theoremstyle{definition}
\newtheorem{definition}[theorem]{Definition}
\newtheorem{assumption}[theorem]{Assumption}
\newtheorem{remark}[theorem]{Remark}
\theoremstyle{remark}
\newtheorem{example}[theorem]{Example}
\newtheorem{question}[theorem]{Question}

\title{The Unified Relational Framework:\\
Geometric Derivation from the Reuleaux Triangle,\\
Residual Stress, Optical Modes, and Cyclic Structure}
\author{Chris Peterson}
\address{Independent researcher}
\keywords{Reuleaux geometry, Force Angle, Angular Shortfall, residual stress, Pressure Shell,
strong nuclear field, pressure length, optical modes, cyclic cosmology}

\begin{document}



\begin{abstract}
We derive a unified geometric framework from the single assumption that the
generating circles of the Reuleaux triangle are retained. This yields the Force
Angle Law, the irreducible Angular Shortfall, and the four-sphere membrane body.
The strong nuclear field is taken as the primordial bonding capacity; its geometric
signature is the permanent shortfall in three-fold closure. Particles are modelled
as layered structures (heart, skin, membrane) carrying a residual-stress dipole
$P/p$ forced by the shortfall. From the same geometry we obtain a pure geometric
pressure length $\ell = s/\sqrt{10}$, elastic moduli ratios, and a five-mode
optical classification of residual-field organisation. Absolute scale is closed by
adopting $\hbar c$ with fundamental length $s\approx 0.84\,\mathrm{fm}$. Dual
dissipation channels (caustic $=$ energetic; solvent boundary $=$ sympathetic) and
aggregation/dissolution pathways at $N_c^*=3$ complete the dynamical picture.
All claims carry explicit status labels.
\end{abstract}

\maketitle

%----------------------------------------------------------------------

\section{Introduction}
%----------------------------------------------------------------------
The behaviour of the strong nuclear field is conventionally described within
Quantum Chromodynamics. Here we adopt a different point of departure: a single
geometric assumption---that the three generating circles of the Reuleaux triangle
are retained as an active part of the construction. From this assumption a
continuous chain follows: the Force Angle Law, the Angular Shortfall, the
four-sphere membrane body, the layered particle, residual stress, helical stacking,
optical organisation of the residual field, and cosmic cyclicity.

The present version substantially extends earlier drafts. In addition to the
geometric skeleton we now include: (i)~the residual-stress dipole and its
combinatorial necessity; (ii)~the pressure length as a pure geometric distortion
scale; (iii)~elastic moduli ratios fixed by the tetrahedron; (iv)~a five-mode
optical taxonomy controlled by the shortfall; (v)~dual dissipation channels;
(vi)~absolute scale closure with $\hbar c$ and a hadronic choice of $s$; and
(vii)~refined aggregation and dissolution pathways at $N_c^*=3$.

\section{Foundational Axioms}
\label{sec:axioms}

\begin{definition}[Strong Nuclear Field]
The \textbf{strong nuclear field} is the primordial bonding capacity of this
universe. It remains the continuous source of all stable structure wherever
distinct centres can approach one another.
\end{definition}

\begin{assumption}[Primordial status]
\assumed\ We take the strong nuclear field as primitive. Its existence and
bonding capacity are not derived from any prior structure within this framework.
\end{assumption}

\begin{theorem}[Angular Shortfall]
\derived\ At every scale where three-fold overlap occurs, closure remains
structurally incomplete. The Force Angle $\theta(d)$ approaches but never reaches
$120^\circ$ for any physically realised separation. This irreducible shortfall
prevents perfect static equilibrium and forces the system to remain cyclic.
\end{theorem}

\begin{proof}
The Force Angle law
\[
\theta(d)=\frac{2\pi}{3}-2\arcsin\Bigl(\frac{d}{2s}\Bigr)
\]
is derived from three circles of equal radius. The limit $\theta\to 120^\circ$ as
$d\to 0$ is never attained for finite $d>0$. At every intersection of any two
circles the third is absent by construction. In-plane closure is an infinite
approach; resolution requires a fourth centre off-plane.
\end{proof}

\section{Reuleaux Geometry and the Membrane Body}
\label{sec:geometry}

\subsection{Force Angle and shortfall}
Three circles of radius $s$ centred at the vertices of an equilateral triangle of
side $d$ produce
\[
\Delta\theta(d):=\frac{2\pi}{3}-\theta(d)=2\arcsin\Bigl(\frac{d}{2s}\Bigr)>0
\quad\text{for all }d>0.
\]
At the tetrahedral scale $d=s$ one has the critical landmark
\[
\Delta\theta_c=\frac{\pi}{3}.
\]

\subsection{Four-sphere membrane}
Lifting to four spheres of radius $s$ at the vertices of a regular tetrahedron of
side $s$ yields a closed membrane shell:
\begin{align*}
r_{\mathrm{inner}}&=\Bigl(1-\frac{\sqrt{6}}{4}\Bigr)s,\\
r_{\mathrm{outer}}&=\Bigl(1+\frac{\sqrt{6}}{4}\Bigr)s,\\
\Delta r_m&=\frac{\sqrt{6}}{2}s.
\end{align*}
The membrane is the region inside some but not all spheres---the three-dimensional
analogue of the two-but-not-three arc condition.

\subsection{Layered particle}
Overlap count partitions the particle into:
\begin{itemize}
\item \textbf{Heart} --- 4-of-4 overlap (maximal bonding; primary SNF locus);
\item \textbf{Skin} --- predominantly 3-of-4;
\item \textbf{Membrane} --- 1-of-4 / 2-of-4 (imperfect, leaky).
\end{itemize}

\section{Residual Stress and the Dipole}
\label{sec:residual}

\begin{definition}[Residual stress]
\derived\ Residual stress is the mechanical expression of the Angular Shortfall:
\[
\sigma_{\mathrm{res}}=\kappa\cdot\Delta\theta,\qquad\kappa\neq 0.
\]
\end{definition}

\begin{proposition}[Residual-stress dipole]
\derived\ Because residual stress must be dynamically active, $\kappa$ cannot
vanish. The only two continuous possibilities are opposite signs:
\[
P>0\quad\text{(tension)},\qquad p<0\quad\text{(pressure)}.
\]
Neither may be zero; both act on the same geometric medium. The unreachable null
separates them but does not forbid imperfect exchange.
\end{proposition}

The Kinetic optical mode (Section~\ref{sec:optics}) is the ordinary living regime
in which both signs remain active---the default state of residual organisation.

\section{Pressure Length and Elastic Moduli}
\label{sec:moduli}

\subsection{Pressure length}
From the Boerdijk--Coxeter rise $H=s/\sqrt{10}$ the critical linear strain is
\[
\varepsilon_c=\frac{1}{\sqrt{10}}.
\]
The \textbf{pressure length} quantifies residual-stress distortion of the cell.
Its magnitude is pure geometry:
\[
\ell=\varepsilon_c\,s=\frac{s}{\sqrt{10}}.
\]
Six combinatorial regimes relate residual pressure to $\ell$:

\begin{center}
\begin{tabular}{lll}
\toprule
Symbol & Condition & Meaning \\
\midrule
$\mathrm{pp}$ & both low & Highest compression \\
$\mathrm{P}\ell$ / $\mathrm{pl}$ & pressure dominates or strong compression & Strong compression \\
$\ell\mathrm{P}$ / $\mathrm{pL}$ & length dominates & Moderate compression / length-led \\
$\mathrm{PL}$ & pressure $=$ length (positive) & Balanced equality (positive) \\
$\mathrm{pl}$ & pressure $=$ length (negative) & Balanced equality (negative) \\
$\mathrm{PP}$ & both high & Highest stretch \\
\bottomrule
\end{tabular}
\end{center}

Primary comparison axis:
\[
\mathrm{P}\ell \;\longleftrightarrow\; \ell\mathrm{P}
\]

Full residual cycle:
\[
\mathrm{pp}\;\to\;\mathrm{pl}\;\to\;\mathrm{PL}\;\to\;\mathrm{pL}\;\to\;\mathrm{PP}\;\to\;\mathrm{Pl}\;\to\;\mathrm{pp}.
\]

\subsection{Elastic moduli from the tetrahedron}
\derived\ (ratios). Identifying residual stress with the load that produces the
geometric critical strains yields pure geometric ratios:
\begin{align*}
Y&=\sqrt{10}\,|\sigma_{\mathrm{res}}|,\\
K&=\frac{\sqrt{10}}{3}\,|\sigma_{\mathrm{res}}|,\\
G&=\frac{3}{\pi}\,|\sigma_{\mathrm{res}}|,
\end{align*}
where $Y$, $K$ and $G$ are Young-type, bulk and shear moduli respectively.

\section{Absolute Scale}
\label{sec:scale}

Pure geometry supplies one length $s$ and pure numbers. Stress, energy and mass
require a dimensionful bridge. We adopt $\hbar c$ as that external constant
\assumed:
\[
|\sigma_{\mathrm{res}}|=C\,\frac{\hbar c}{s^2},
\]
with $C$ a pure geometric prefactor of order $10$--$50$.

Reading residual-stress energy of the membrane volume as rest energy
\interpretive\ gives
\[
m\sim\frac{\hbar}{c\,s}.
\]
The fundamental length is fixed at the proton charge radius
\interpretive:
\[
s\approx 0.84\,\mathrm{fm}.
\]
With $\hbar c\approx 197.3\,\mathrm{MeV\cdot fm}$ one obtains
\[
\frac{\hbar c}{s^2}\approx 280\,\mathrm{MeV\cdot fm^{-1}},\qquad
m\sim 235\,\mathrm{MeV}/c^2
\]
(order of light-hadron masses), and
\[
\ell\approx 0.27\,\mathrm{fm}.
\]

\begin{remark}
Identifying $s$ with the Planck length would close the framework onto Planck units
but place the particle at the Planck mass. The hadronic choice is preferred if the
layered cell is to describe ordinary matter.
\end{remark}

\section{Generalised Optics}
\label{sec:optics}

All optical organisation of the residual electromagnetic field is controlled by
the Angular Shortfall. Residual stress is the deflecting medium; no separate
optical axiom is required.

\subsection{Propagation zones}
\begin{center}
\begin{tabular}{lll}
\toprule
Zone & Overlap & Character \\
\midrule
Heart & 4-of-4 & Nearest vacuum; speed $\approx c$; minimal deflection \\
Skin & 3-of-4 & Intermediate residual; mild deflection \\
Membrane & 1--2 of 4 & Strongest residual; leakage, focusing, diffusion \\
Exterior & 0 & Ambient \\
\bottomrule
\end{tabular}
\end{center}

\subsection{Five optical modes}
\conceptual\ The residual field organises into five geometric regimes:

\begin{center}
\begin{tabular}{lll}
\toprule
Mode & Character & Geometric seat \\
\midrule
Caustic & Hyper-concentrated focus & Envelope of rays folded by shortfall gradients \\
Prismatic & Dispersion into components & Differential deflection at incomplete loci \\
Kinetic & Living intermediate residual & Ordinary residual motion; seat of $P/p$ \\
Scotopic & Dim, reduced intensity & Long membrane paths \\
Cryptic & Flat, outline-erasing diffusion & Solvent-boundary smoothing \\
\bottomrule
\end{tabular}
\end{center}

\subsection{Dual dissipation channels}
\begin{itemize}
\item \textbf{Caustic $=$ energetic}: local concentration, irreversible conversion
at singularity approach.
\item \textbf{Solvent boundary $=$ sympathetic}: resonant transfer across the
membrane between opposing residual domains.
\end{itemize}

\subsection{Refractive dipole}
Parallel to residual stress,
\[
N>1\quad(c_{\mathrm{med}}<c),\qquad n<1\quad(\text{phase }c_{\mathrm{med}}>c).
\]
Both are combinatorially forced once residual response is non-zero.

\section{Dynamics near the Null}
\label{sec:dynamics}

As a local configuration approaches a $120^\circ$ null:
\begin{enumerate}
\item shortfall-driven deflection intensifies;
\item residual frequency rises (Euler-disk character);
\item caustic formation becomes more probable;
\item at a critical threshold an equivariant reorganisation
(compatible with $\mathbb{Z}_3$ unfolding) can occur and residual-stress sign may flip;
\item both dissipation channels are transiently enhanced.
\end{enumerate}
\conceptual\ The chords of the envelope construction, once idealisation is dropped,
are candidates for interaction lines of the residual web; their caustic is the
collective signature of that imperfect web.

\section{Aggregation and Dissolution}
\label{sec:pathways}

\begin{definition}[Critical contact number]
\derived\ $N_c^*=3$. A regular tetrahedron requires four centres; three contacts
are necessary and sufficient for one particle to complete a new tetrahedral unit.
The same count attaches a new tetrahedron to an exposed face of a Boerdijk--Coxeter
helix.
\end{definition}

\begin{itemize}
\item If $N_c<3$ the particle follows the \textbf{dissolution pathway}: ordered
layering is lost and capacity returns to the ambient field.
\item If $N_c\ge 3$ the particle follows the \textbf{aggregation pathway}: residual
tension couples contacting membranes into larger composite structure.
\end{itemize}

Nucleation of a new cell is geometrically natural at a charged opposing ($P$ against
$p$) interface when residual accumulation near critical shortfall discharges into
the solvent gap. Dissolution is open for under-coordinated particles ($N_c<3$),
with elevated rate on the periphery and under solitary caustic spikes.

\section{Helical Dynamics}
\label{sec:helix}

Face-to-face aggregation produces the Boerdijk--Coxeter helix with parameters fixed
by the tetrahedron:
\[
\theta_{BC}=\arccos(-2/3),\qquad H=\frac{s}{\sqrt{10}},\qquad R=\frac{3s\sqrt{3}}{10}.
\]
Nine-stack assemblies exhibit frequency locking driven by shared residual tension.
Chiral flow and radial breathing of outer membranes are geometric consequences of
the permanent shortfall. Under $S^3$ topology the helix closes after thirty
tetrahedra as the $(2,3)$ torus knot (trefoil).

\section{Cosmological Implications}
\label{sec:cosmology}

Cyclicity is forced by the Angular Shortfall: no configuration reaches perfect
static equilibrium. Progressive flattening, saturation, residual release and
renewed three-dimensional propagation constitute a complete geometric cycle.
Hierarchical continuity means the same rules operate from the single particle to
the largest assemblies. Quantitative cosmological predictions remain \openq.

\section{Computational Realisations}
\label{sec:computational}

Interactive realisations implement the derived geometry without free parameters
beyond $s$:
\begin{itemize}
\item layered particle with residual-stress cycle, pressure-length breathing, and
optical-mode tagging;
\item stacked particles with solvent boundaries and sympathetic coupling;
\item nucleation field with interface-charge accumulation on opposing residual
contacts, aggregation at $N_c^*\ge 3$, and dissolution for $N_c<3$.
\end{itemize}
These confirm internal consistency of the geometric chain and the qualitative
dynamics of residual stress, nucleation and dissolution.

\section{Conclusions}
\label{sec:conclusions}

\subsection{Established}
From retained generating circles: Angular Shortfall, membrane body, layered
particle, residual-stress dipole, pressure length, elastic moduli ratios, optical
modes and dual dissipation, contact rule $N_c^*=3$, helical parameters, and the
necessity of cyclicity.

\subsection{Status of main claims}
\begin{center}
\begin{tabular}{ll}
\toprule
Claim & Status \\
\midrule
Angular Shortfall as geometric necessity & Derived \\
Layered particle from four-sphere overlap & Derived \\
Residual-stress dipole $P/p$ & Derived \\
Pressure length $\ell=s/\sqrt{10}$; regimes $\mathrm{Pl},\mathrm{pL},\mathrm{PL},\mathrm{pl}$ & Derived \\
Elastic moduli ratios $Y,K,G$ & Derived (ratios) \\
Five optical modes & Conceptual (geometry-seated) \\
Dual dissipation channels & Conceptual \\
Scale closure with $\hbar c$, $s\approx 0.84\,\mathrm{fm}$ & Assumed / interpretive \\
$N_c^*=3$ pathways & Derived \\
Cyclicity forced by shortfall & Derived \\
EM field from membrane leakage & Motivated hypothesis \\
Quantitative cosmology & Open \\
\bottomrule
\end{tabular}
\end{center}

\subsection{Open questions}
\begin{enumerate}
\item Precise geometric prefactor $C$ in $|\sigma_{\mathrm{res}}|=C\,\hbar c/s^2$.
\item Full 3D ray-trace with shortfall-controlled deflection.
\item Detailed laboratory mapping of $P/p$ and $N/n$.
\item Whether $s$ is exactly the proton radius or a geometric multiple thereof.
\item Quantitative cosmological confrontation.
\end{enumerate}

\subsection{Final remark}
The Angular Shortfall is small, permanent, and present at every scale. From that
single fact the layered particle, residual stress, optical organisation, helical
structure and the necessity of cycles all follow. The universe described here is
never at rest, not because of an external driver, but because perfect closure was
never available.

\appendix

\section{Glossary of Terms}
\label{app:glossary}

\begin{description}
\item[Reuleaux triangle]
Intersection of three circles of equal radius $s$ centred at the vertices of an
equilateral triangle. Constant width $s$.

\item[Generating circles]
The three full circles of radius $s$ retained as geometrically real throughout
the construction (not discarded after the boundary arcs are drawn). This is the
foundational assumption of the framework.

\item[Force Angle]
The opening angle formed at a vertex by any two generating circles:
\[
\theta(d)=\frac{2\pi}{3}-2\arcsin\Bigl(\frac{d}{2s}\Bigr),\qquad d\in[0,s\sqrt{3}].
\]
Approaches but never reaches $120^\circ$ for finite $d>0$. Formerly called the
Fork Angle; renamed to emphasise its role as the seat of residual mechanical force.

\item[Angular Shortfall (Energetic loss)]
The permanent angular deficit from perfect three-fold closure:
\[
\Delta\theta(d)=2\arcsin\Bigl(\frac{d}{2s}\Bigr)>0\quad\text{for all }d>0.
\]
Critical value at the tetrahedral edge ($d=s$): $\Delta\theta_c=\pi/3$.
Root cause of residual stress, residual angular momentum, and cyclicity.

\item[Pressure Shell]
The four-sphere shell formed when three circles are lifted to four spheres of
radius $s$ at the vertices of a regular tetrahedron of side $s$. Inner radius
$(1-\sqrt{6}/4)s$, outer radius $(1+\sqrt{6}/4)s$, thickness $(\sqrt{6}/2)s$.
Formerly called the membrane body.

\item[Extended nucleus]
The layered composite object formed by the Pressure Shell:
\begin{itemize}
\item \textbf{Heart} --- maximal 4-of-4 overlap (core, strongest bonding);
\item \textbf{Skin} --- predominantly 3-of-4 overlap (transitional);
\item \textbf{Pressure Shell} --- 1-of-4 or 2-of-4 overlap (imperfect, leaky outer layer).
\end{itemize}
Formerly called the layered particle.

\item[Residual stress]
Mechanical expression of the Angular Shortfall:
\[
\sigma_{\mathrm{res}}=\kappa\cdot\Delta\theta,\qquad\kappa\neq 0.
\]

\item[Residual-stress dipole]
The two opposite signs forced by the shortfall:
$P>0$ (tension / stretch) and $p<0$ (pressure / compression).
Neither sign may vanish while the shortfall exists.

\item[Pressure length $\ell$]
Pure geometric distortion scale (no stress symbol required):
\[
\ell=\frac{s}{\sqrt{10}}\quad\text{(pressure length; pure geometry)}.
\]
Derived from critical strain $\varepsilon_c=1/\sqrt{10}$. Distinct from residual stress $\sigma_{\mathrm{res}}$.

\item[Pressure--length comparison]
Primary axis $\mathrm{P}\ell \leftrightarrow \ell\mathrm{P}$ (pressure greater than length $\leftrightarrow$ length greater than pressure).

\item[Pressure-length regimes]
Combinatorial states of residual pressure relative to $\ell$:
$\mathrm{pp}$ (highest compression), $\mathrm{pl}$ / $\mathrm{P}\ell$ (strong compression),
$\mathrm{pL}$ / $\ell\mathrm{P}$ (length-led), $\mathrm{PL}$ (balanced positive),
$\mathrm{pl}$ (balanced negative), $\mathrm{PP}$ (highest stretch), $\mathrm{Pl}$ (moderate stretch).

\item[Elastic moduli]
Young-type, bulk and shear ratios fixed by tetrahedral geometry:
\[
Y=\sqrt{10}\,|\sigma_{\mathrm{res}}|,\quad
K=\frac{\sqrt{10}}{3}\,|\sigma_{\mathrm{res}}|,\quad
G=\frac{3}{\pi}\,|\sigma_{\mathrm{res}}|.
\]

\item[Inertial mass]
\[
m\sim\frac{\hbar}{c\,s}.
\]
With $s\approx 0.84\,\mathrm{fm}$ this yields hadronic-scale mass.

\item[Residual angular momentum]
\[
\mathbf{L}_{\mathrm{res}}\propto\kappa\cdot\Delta\theta\times\mathbf{r}.
\]
The Angular Shortfall acting at a lever arm produces permanent rotational bias
(honeycombing, vortex action, Pressure Shell spin).

\item[Critical contact number]
\[
N_c^*=3.
\]
Minimum membrane--membrane contacts for aggregation; below this the dissolution
pathway is open.

\item[Aggregation pathway]
$N_c\ge 3$: residual tension couples contacting membranes into larger composite
structure (including Boerdijk--Coxeter helical growth).

\item[Dissolution pathway]
$N_c<3$: ordered layering is lost; capacity returns to the ambient field.

\item[Five optical modes]
Geometric regimes of residual electromagnetic organisation controlled by the
shortfall: Caustic (hyper-concentrated focus), Prismatic (dispersion),
Kinetic (living intermediate residual; seat of $P/p$), Scotopic (dim),
Cryptic (flat diffusion).

\item[Dual dissipation]
Caustic channel $=$ energetic; solvent boundary $=$ sympathetic transfer across
opposing residual domains.

\item[Boerdijk--Coxeter helix]
Face-to-face stacking of tetrahedra:
\[
\theta_{BC}=\arccos(-2/3),\quad H=\frac{s}{\sqrt{10}},\quad R=\frac{3s\sqrt{3}}{10}.
\]

\item[Cyclicity]
Forced by the Angular Shortfall: no configuration reaches perfect static
equilibrium. The system evolves through flattening, saturation, residual release
and renewed propagation.

\item[Strong nuclear field]
Primordial bonding capacity of the universe; taken as primitive. Its geometric
signature is the permanent shortfall in three-fold closure.
\end{description}

\section{Complete Geometric Derivations}
\label{app:derivations}

\subsection{A.1 Force Angle}
\[
\theta(d)=\frac{2\pi}{3}-2\arcsin\Bigl(\frac{d}{2s}\Bigr),\qquad d\in[0,s\sqrt{3}].
\]

\subsection{A.2 Angular shortfall}
\[
\Delta\theta(d)=2\arcsin\Bigl(\frac{d}{2s}\Bigr)>0\quad\text{for all }d>0;
\qquad\Delta\theta(s)=\frac{\pi}{3}.
\]

\subsection{A.3 Membrane radii}
\[
r_{\mathrm{inner}}=\Bigl(1-\frac{\sqrt{6}}{4}\Bigr)s,\quad
r_{\mathrm{outer}}=\Bigl(1+\frac{\sqrt{6}}{4}\Bigr)s,\quad
\Delta r_m=\frac{\sqrt{6}}{2}s.
\]

\subsection{A.4 Pressure length and moduli}
\[
\varepsilon_c=\frac{1}{\sqrt{10}},\quad
\ell=\frac{s}{\sqrt{10}}\quad\text{(pressure length; pure geometry)}.
\]
Regimes of pressure relative to length:
\[
\mathrm{pp},\;\mathrm{pl},\;\mathrm{pL},\;\mathrm{PL},\;\mathrm{PP},\;\mathrm{Pl}
\]
(see Glossary, Appendix~\ref{app:glossary}).
Elastic moduli:
\[
Y=\sqrt{10}\,|\sigma_{\mathrm{res}}|,\quad
K=\frac{\sqrt{10}}{3}\,|\sigma_{\mathrm{res}}|,\quad
G=\frac{3}{\pi}\,|\sigma_{\mathrm{res}}|.
\]

\subsection{A.5 Scale}
\[
|\sigma_{\mathrm{res}}|=C\frac{\hbar c}{s^2},\qquad
m\sim\frac{\hbar}{c s},\qquad
s\approx 0.84\,\mathrm{fm}.
\]

\subsection{A.6 Boerdijk--Coxeter}
\[
\theta_{BC}=\arccos(-2/3),\quad H=\frac{s}{\sqrt{10}},\quad R=\frac{3s\sqrt{3}}{10}.
\]

\subsection{A.7 Critical contact}
\[
N_c^*=3.
\]

\begin{thebibliography}{99}
\bibitem{reuleaux1876} Reuleaux, F. (1876). \emph{Theoretische Kinematik}. Vieweg.
\end{thebibliography}

\end{document}
